\section{Introduction}
\label{sec:intro}

Postoperative atrial fibrillation (POAF) affects roughly one in four patients undergoing isolated coronary artery bypass grafting (CABG) and carries substantial risk of perioperative stroke, prolonged intensive care stay, hospital readmission, higher expenditure, and reduced long-term survival.\cite{aranki1996predictors,villareal2004,dobrev2019postoperative,caldonazo2023meta,gaudino2023ehj} Targeted preventive interventions---including posterior left pericardiotomy,\cite{gaudino2021palacs} prophylactic amiodarone or colchicine, and intensified postoperative telemetry\cite{chatterjee2026sts}---offer clinical benefit, yet universal prophylaxis exposes lower-risk surgical patients to avoidable adverse effects. Accurate, well-calibrated risk stratification available before skin incision is therefore critical both for individualized surgical prophylaxis and for enriching clinical trials.

Despite rapid growth in clinical machine learning, published POAF models report areas under the receiver-operating-characteristic curve (ROC-AUCs) ranging from 0.58 to 0.98.\cite{fields2025poaf,lopeztuesta2026poaf} Two recent PROBAST systematic reviews demonstrated that nearly all published algorithms carry high risk of bias and suffer severe external shrinkage, falling from a median apparent AUC of 0.71 to a median external AUC of 0.61.\cite{fields2025poaf,lopeztuesta2026poaf} Apparent AUCs exceeding 0.85 almost always reflect intraoperative or intensive-care predictors,\cite{mathew2004multicenter,magee2007atrial,ma2025noaf,cao2025mimic4} mixed valve-and-CABG surgical cohorts,\cite{mariscalco2014poaf,burgos2021comaf} single-center internal splits,\cite{cho2026poaf,shakhgeldyan2026cabg} continuous waveform or genomic modalities outside routine registries,\cite{han2025aiecg,khalpey2026digital} or differences between registry-coded treatment-requiring AF and continuous telemetry ascertainment.\cite{filardo2017telemetry,perezgrovas2023jaha}

Compounding these comparability issues, national surgical registries expose two deeper methodological vulnerabilities. First, naive outcome filtering and default final-year temporal splits distort external evaluation: in the Society of Thoracic Surgeons Adult Cardiac Surgery Database (STS ACSD),\cite{sts2014v281,shahian2018sts,obrien2018sts2} conditional skip logic leaves uncomplicated recoveries blank, early table columns conceal postoperative durations, and calendar-2017 holdouts retain 99.5\% hospital overlap with development while mixing half-year coverage with missingness drift. Second, conventional binary $\{0,1\}$ supervision compels models to fit noisy outcomes heavily influenced by unmeasured intraoperative perturbations,\cite{echahidi2008mechanisms,dobrev2019postoperative} leaving gradient-boosted trees and frontier tabular foundation models (TabPFN-3.5,\cite{hollmann2025tabpfn,jaeger2026tabpfn35} KumoTabular-Large v2,\cite{qu2026kumotabular} and TabFM v1.0.0\cite{kong2026tabfm}) clustered in a narrow performance band.

Here we audit all 368,132 isolated-CABG procedures across 1,091 US surgical centers in STS ACSD v2.81 and evaluate a dual-paradigm modeling framework pairing 1,000-generation evolutionary continuous-target learning (\textbf{EvoTarget-FT}) with re-engineered GPU-guided parallel genetic programming (\textbf{JAIN Ultra GP v2}). We benchmark our approach against six established bedside clinical risk scores, six world-leading binary tabular learning families, and a 7$\times$-larger national full-registry reference under both strictly unseen-hospital external transport (0.0\% hospital overlap across 74 centers) and a controlled in-network seen-hospital deployment regime in which all 1,091 hospitals appear in training, evaluating discrimination, calibration, Decision Curve net benefit, demographic equity, and institutional site recalibration.
